\clearpage\NoSourceHead\SourceLeaf{389}
\begingroup
\makeatletter\tagsleft@true\let\veqno\@@leqno\makeatother
\HistoricTitle{[ZUR LEHRE VON DEN REIHEN.]}
\SourceGroup{I}
\HistoricTitle{NEUE METHODE DIE SUMME DER DIVERGIRENDEN REIHE\\
\(1-1+2-6+24-\text{etc.}\ [=]\ 0{,}5963\,{.}\,{.}\,{.}\)\\
ZU FINDEN.}
\HistoricCitation{[Ein Blatt in Fa, Kapsel 46a]}
\begin{center}1.\end{center}
\noindent\hspace*{\parindent}Es sei
\begin{equation}
x-x^2+2x^3-6x^4+\text{etc.}=\xi\tag*{[1]}
\end{equation}
\noindent und man wird finden
\begin{equation}
\xi+xx\xi'-x=0.\tag*{[2]}
\end{equation}
\noindent Diese Different[ial] Gl[eichung] muss so integrirt werden, dass für \(x=0,\ \xi=0\) werde; der Werth von \(\xi\), für \(x=1\), den wir \(c\) nennen[\textsuperscript{*})], ist die gesuchte Summe.
\begin{center}2.\end{center}
\noindent\hspace*{\parindent}Wir bekommen durch fortgesetzte Differentiation folgende Gleich[ungen]:
\[
\begin{gathered}
\xi+xx\xi'-x=0\\
(1+2x)\xi'+xx\xi''-1=0\\
2\xi'+(1+4x)\xi''+xx\xi'''=0\\
6\xi''+(1+6x)\xi'''+xx\xi^{\mathrm{IV}}=0\\
12\xi'''+(1+8x)\xi^{\mathrm{IV}}+xx\xi^{\mathrm{V}}=0\\
\text{etc.}
\end{gathered}
\]
\par\medskip\noindent\rule{.24\linewidth}{.4pt}\par\smallskip
{\small\noindent\hspace*{\parindent}[\textsuperscript{*}) In der Handschrift wird dieser Wert mit \(e\) bezeichnet.]\par}
\clearpage\endgroup
